\chapter*{\textarabic{ملخص}}
\addcontentsline{toc}{chapter}{ملخص}

\begin{Arabic}

يتناول مشروع نهاية الدراسة هذا تقليص مساحة الهجوم في سياق إدارة الأنظمة الحساسة، وذلك من خلال بنية تحتية متكاملة قائمة على مبدأ "انعدام الثقة" (Zero Trust). تكمن الإشكالية الرئيسية التي تم تحديدها في وجود صلاحيات وصول امتيازية دائمة، تفتقر إلى التتبع الدقيق ولا ترتبط بشكل كافٍ بمستوى حساسية الأصول الرقمية.

تعتمد الحلول المقترحة على دمج خمس ركائز أساسية: Authentik لإدارة الهويات (IAM)، و Tailscale ACL للتحكم في الوصول إلى الشبكة، وتطبيق ويب Flask مع قاعدة بيانات PostgreSQL لحوكمة الطلبات، ومحرك Python (Just-in-Time) للأتمتة والموافقة أو الإلغاء التلقائي، بالإضافة إلى Wazuh لضمان قابلية المراقبة والتتبع الأمني.

يفرض مسار العمل المستهدف عدم وجود أي وصول امتيازي دائم؛ حيث يتم طلب كل رفع للصلاحيات وتبريره والموافقة عليه، مع تقييده بفترة زمنية محددة (TTL) ثم سحبه تلقائيًا. تشمل سيناريوهات التحقق: المصادقة القوية، طلب الوصول الآني (JIT)، تفعيل الـ ACL، والوصول عبر بروتوكول SSH خلال النافذة الزمنية المسموح بها، متبوعًا برفض الوصول فور انتهاء الصلاحية.

تُظهر النتائج المتوقعة تحسنًا ملحوظًا في الوضع الأمني من خلال: إلغاء الوصول الدائم غير المبرر، التتبع الكامل لدورة حياة الوصول، وتعزيز القدرة على التدقيق والامتثال للمعايير الأمنية.

\textbf{الكلمات المفتاحية:} Zero Trust، IAM، PIM، PAM، Authentik، Tailscale، JIT، ACL، Wazuh، الأمن السيبراني.

\end{Arabic}

\clearpage